\subsection{鏡共鳴を二曲線へ還元し、一曲面だけを残す}
\begin{proposition}
鏡共鳴 $t=(\rho^2-1)/2$ のH3候補は
\[
 J_m=4\rho^2s^2+(\rho+1)^2(\rho^2+1)^2=0
\]
に限られる。
従って、両top非零のmixed4 $(2,2)$ 型に残るH3候補は、主chartの一曲面 $L=0$ に全て含まれる。
これはその曲面の許容条件や全H3被覆の完成を意味しない。
\end{proposition}
\begin{proof}
$a=0$ は排除済みなので
\[
 C=k[\rho,s,U_m^{-1}],\qquad
 U_m=\rho(\rho^2-1)\bigl(4s^2+(\rho+1)^2(\rho^2+3)\bigr)
\]
上で考える。これは鏡共鳴における正確な必要局所化である。
前節と同じ正三次関係行列の核 $K$ と、正常行加群 $F\subset K$ について
\[
 B_mJ_m K\subset F,\qquad B_m=4s^2+(\rho-1)(\rho+1)^3
\]
を標数0の局所化加群計算で直接検証する。
AS4最大rank開集合上で $K$ は正三次商の全双対へ底変更するため、
正常三次が非零ならば $B_mJ_m=0$ でなければならない。

$B_m=0$ の全域を調べる。
\[
 \rho=\frac{1-w^2}{1+w^2},\qquad
 s=\frac{2w}{(1+w^2)^2},\qquad
 t=-\frac{2w^2}{(1+w^2)^2},\qquad
 a=\frac{4w^2}{(1+w^2)^2}
\]
と置く。必要条件は $w(w^2-1)(w^2+1)\ne0$ である。
逆写像 $w=2s/(\rho+1)^2$ があるので、除外済みの $\rho=\pm1,0$ を除く全曲線と両符号を覆う。

この曲線に対する多項式四次汎関数を入力
\[
 \texttt{cubic\_order2\_22\_secant\_mirror\_B\_input.json}
\]
に与える。元の全288関係を消す恒等式を $\mathbb Q[w]$ 上で再生し、
全正自由三次32語への正常行列の10minorを調べる。
そのminorは非零で、零点は
\[
 w(w^2-1)(w^2+1)(w^4-w^2+1)=0
\]
に限られる。
従って $w^4-w^2+1\ne0$ なら正常核は零となる。
残る点も恒等式
\[
 J_m=\frac{16(w^4-w^2+1)}{(1+w^2)^4}
\]
により $J_m=0$ に含まれる。
これで鏡共鳴の全候補を $J_m=0$ へ縮めた。

最後に、元の主chartでの多項式 $L$ は
\[
 L\bigm|_{t=(\rho^2-1)/2}=\frac{\rho-1}{2}J_m
\]
を満たす。
前節までの全座標曲面排除と主chart開集合排除を合わせれば、
未解決の全候補が一曲面 $L=0$ に含まれることが従う。

\texttt{verify\_cubic\_order2\_22\_secant\_mirror\_support.py}
は局所化加群の所属、曲線上の元関係汎関数、正常10minor及び二つの有理恒等式を再生する。
$J_m=0$ 上の正常切断を全て排除したとは主張しない。
\end{proof}
